\subsection{Recovering individual factors}

Let \(B_j=\prod_{i\le j} b_i\), \(B_0=1\), and let \(E_{j,h}\) be
the \(j\)-th elementary symmetric sum of the sequence \((b_i^h)_i\)
for integer \(h\ge1\). In particular \(E_{0,h}=1\). The logarithm
of \eqref{eq:all-real-product} near zero and Newton's identities give
\[
 p_n=-n[v^n]\log G(v)=\sum_i b_i^n,qquad
 E_{j,h}=\frac1j\sum_{n=1}^j(-1)^{n-1}p_{nh} E_{j-n,h} .
\]
The recurrence follows first for a finite head and then by absolute
convergence. Each \(E_{j,h}\) uses even moments through order \(2jh\).
Since \(B_j\) is the largest indexed product in \(E_{j,1}\),
\(B_j^h\le E_{j,h}\le B_j^{h-1}E_{j,1}\).
For \(h>1\) define
\(\underline B_{j,h}=(E_{j,h}/E_{j,1})^{1/(h-1)}\) and
\(\overline B_{j,h}=E_{j,h}^{1/h}\), with both equal to one at \(j=0\).
Then
\begin{equation}\label{eq:factor-interval}
 \sqrt{\frac{\underline B_{k-1,h}}{\overline B_{k,h}}}
 \le x_k\le
 \sqrt{\frac{\overline B_{k-1,h}}{\underline B_{k,h}}}.
\end{equation}
In particular \(x_k=\lim_{h\to\infty}(E_{k-1,h}/E_{k,h})^{1/(2h)}\)
for each fixed rank, including multiplicities. For finite input
intervals the alternating recurrence requires certified coefficient
errors before dividing by positive lower endpoints.

If the preceding factors are known, a separate estimate uses
\(s_h^{(k)}=p_h-\sum_{i<k}b_i^h\). The ratio
\(s_{h+1}^{(k)}/s_h^{(k)}\) is a weighted average of the factors
from index \(k\) onward. Subtract the average divided by \(b_k\)
from one, then use \(s_h^{(k)}\ge b_k^h\) and bound each
\((b_i/b_k)^{h-1}\) for \(i>k\) by the first such term. This gives
\[
 0\le 1-\frac{s_{h+1}^{(k)}}{b_k s_h^{(k)}}
 \le (b_{k+1}/b_k)^{h-1}\sum_{i>k} b_i/b_k .
\]
Strict separation gives geometric decay after the tail is bounded.
The subtraction also requires precision on the preceding factors
\cite[\S3]{Founding07}.

\paragraph{Exchanges across one rank.}
For \(b_{j+1}\le\theta_j\le b_j\), write
\[
 A_j=\sum_{i\le j}(\theta_j/b_i)^h,qquad
 C_j=\sum_{i>j}(b_i/\theta_j)^h .
\]
A term in \(E_{j,h}/B_j^h\) removes and inserts the same number of
indices across the cut at \(j\). All powers of \(\theta_j\)
therefore cancel. Writing \(e_r(\cdot)\) for an elementary symmetric
function here gives
\begin{equation}\label{eq:exchange-bound}
\begin{aligned}
 1\le \frac{E_{j,h}}{B_j^h}
 &=\sum_{r=0}^j e_r(((\theta_j/b_i)^h)_{i\le j})
                  e_r(((b_i/\theta_j)^h)_{i>j})\\
 &\le\sum_{r=0}^{\infty}\frac{(A_jC_j)^r}{(r!)^2}
 \le e^{A_jC_j}.
\end{aligned}
\end{equation}
The symmetric sums are bounded by the indicated powers divided by
\(r!\) on each side, by expanding the powers and counting the terms
with distinct indices. The tail sums converge by \(\sum_i b_i<\infty\).

If \(D_j\ge A_j C_j\) and \(0<L_j\le E_{j,h}\le U_j\), then
\begin{equation}\label{eq:exchange-interval}
 \left(\frac{L_{k-1}e^{-D_{k-1}}}{U_k}\right)^{1/(2h)}
 \le x_k\le
 \left(\frac{U_{k-1}e^{D_k}}{L_k}\right)^{1/(2h)} .
\end{equation}
Use \(D_0=0\) and \(L_0=U_0=1\).
Indeed, \(B_j^h\) lies between \(L_j e^{-D_j}\) and \(U_j\), and
\(x_k^{2h}=B_{k-1}^h/B_k^h\). The logarithm of the ratio of its two
endpoints is
\[
 \frac{D_{k-1}+D_k+\log(U_k/L_k)+\log(U_{k-1}/L_{k-1})}{2h}.
\]
A finite counted head through \(J\ge j+1\) supplies \(D_j\) if a
threshold can be chosen between the upper bound for \(b_{j+1}\) and
the lower bound for \(b_j\). Bound the known weights by interval
endpoints. For the rest, if \(b_{J+1}\le\beta\le\theta_j\) and
\(\sum_{i>J}b_i\le S\), the remaining contribution to \(C_j\)
is at most \((\beta/\theta_j)^{h-1}S/\theta_j\). The upper
endpoint for \(b_J\) can supply \(\beta\), and
\(\sum_i b_i=e_1\) supplies \(S\) after subtracting a lower
bound on the head sum. If both nearest threshold ratios are at
most \(q<1\), then \(A_j\le jq^h\) and
\(C_j\le q^{h-1}\sum_{i>j}b_i/\theta_j\). This gives geometric
decay for fixed \(j\). At a repeated factor \eqref{eq:factor-interval}
still applies.

At high ranks the spacing gives a uniform allowance for each fixed
\(T\). Take \(\theta_j=(b_j b_{j+1})^{1/2}\).
For \(\lfloor j/2\rfloor\le i\le 2j\), the adjacent differences
\(\log b_i-\log b_{i+1}=2\log(1+\Delta_i/x_i)\) are at least
\(c_T/j\) by \eqref{eq:high-rank-geometry}. The nearest weights on
each side of the cut thus cost half a step, and each subsequent
weight at least one additional step. On this range either sum is at
most \(e^{-c_T h/(2j)}/(1-e^{-c_T h/j})\).
On the far head, each base \(\theta_j/b_i\) is at most a fixed
\(q_0<1\) since \(x_{\lfloor j/2\rfloor}/x_j\to1/2\).
On the far tail each base \(b_i/\theta_j\) is also at most
\(q_0<1\) after increasing \(q_0\), since \(x_{2j}/x_j\to2\).
The head costs at most \(j q_0^h\), and the tail costs at most
\(q_0^{h-1}\theta_j^{-1}\sum_{i>2j}b_i=O_T(jq_0^{h-1})\).
For \(h\ge j\) these far bounds and both near sums are
\(O_T(e^{-c'_T h/j})\) with some \(c'_T>0\). For example,
one can absorb the factor \(j\) into half the exponential decay
in \(h\). We can therefore take \(D_j=C_Te^{-c_T'' h/j}\)
in \eqref{eq:exchange-interval} for sufficiently large \(j\)
and all \(h\ge j\).

\paragraph{Uncertainty in the symmetric coefficients.}
For \(\omega=e^{2\pi\ii/h}\), multiplying the rotated products gives
\[
 \prod_{l=0}^{h-1}G(\omega^l v)=\prod_i(1-b_i^h v^h),
 \qquad
 (-1)^j E_{j,h}=[v^{jh}]\prod_{l=0}^{h-1}G(\omega^l v).
\]
The product is locally absolutely convergent, and for \(j\le k\)
only the coefficients through \(kh\) are used. Equip these
truncated polynomials with
\(\|P\|_R=\sum_m |[v^m]P|R^m\), for \(R>0\).
The norm is submultiplicative and rotations are isometries.
Let \(\widetilde G\) be the center polynomial and \(d_R\) the
weighted sum of its coefficient radii in this norm. Put
\(U_R=\|\widetilde G\|_R+d_R\). Replacing the rotated factors
one at a time gives a telescoping sum. Taking its norm and then
extracting degree \(jh\) gives
\begin{equation}\label{eq:exchange-coefficient-error}
 |\widetilde E_{j,h}-E_{j,h}|
 \le h U_R^{h-1}d_R/R^{jh}.
\end{equation}
Enclose any arithmetic error separately.

Write \(M_R=G(-R)=F(\ii\sqrt R)\). The exact truncated norm is
at most \(M_R\). If \(d_R\le\eta M_R\), then
\(U_R\le(1+2\eta)M_R\). Since \(E_{j,h}\ge B_j^h\),
the relative error in \eqref{eq:exchange-coefficient-error} is at most
\[
 h\eta(1+2\eta)^{h-1}
 \left(\frac{M_R}{R^j B_j}\right)^h.
\]
Take \(R=(4\pi k/\log k)^2\). For \(j=k-1,k\),
\begin{equation}\label{eq:exchange-coefficient-growth}
 \frac{M_R}{R^j B_j}
 =\prod_{i\le j}(1+x_i^2/R)\prod_{i>j}(1+R/x_i^2)
 \le e^{C_T k}.
\end{equation}
Every head factor is bounded by a constant since \(x_j^2\asymp R\).
The tail logarithm is at most \(R\sum_{i>j}b_i=O_T(k)\).
Thus coefficient radii \(2^{-p} e_m\) through degree \(kh\)
are sufficient for relative output errors \(O(\tau)\) if
\(p\ge C_T'kh+\log_2(h/\tau)+C_T''\), with constants large enough
to absorb the factor \((1+2\eta)^{h-1}\) and the passage from
radii to endpoints.

\begin{theorem}[Sufficient product input budget]\label{thm:high-rank-exchange}
For each fixed positive all-real \(T\), at sufficiently large rank
\(k\), \(h\ge k\) and relative coefficient enclosure radii at most
\(\tau\le1/4\) for \(E_{k-1,h},E_{k,h}\), the width of
\eqref{eq:exchange-interval} in units of \(\Delta_k\) is
\[
 O_T\!\left(\frac{k}{h}(e^{-c_T h/k}+\tau)\right).
\]
In particular, \(h\asymp_T k(1+\log(1/\varepsilon))\) with a
sufficiently large constant and \(\tau\asymp\varepsilon\)
give width \(O_T(\varepsilon\Delta_k)\). The sufficient even
moment order is \(2kh\), and the sufficient input precision per
normalized coefficient is \(O_T(kh+\log(h/\tau))\) bits.
\end{theorem}
\begin{proof}
Substitute the high-rank bound for \(D_j\) into the logarithmic
width of \eqref{eq:exchange-interval}. Each log coefficient
ratio is \(O(\tau)\). The logarithmic width is \(O_T(1/h)\),
so exponentiating and using \(x_k/\Delta_k\asymp k\) proves
the width estimate. The precision bound follows from
\eqref{eq:exchange-coefficient-error}--\eqref{eq:exchange-coefficient-growth}.
\end{proof}
The sufficient asymptotic constants come from the spacing theorem.
At a specified rank one can use the certified head and tail
allowance above. The symmetric coefficients themselves are evaluated
directly from the moments.
